\section{The exact joint odd CM image on the primary kernel}
\label{sec:exact-joint-odd-cm-image}

We use the groups
\[
 M=\ker\lambda,\qquad
 P_D=D\cap\bigcap_p\ker\Phi_p,\qquad
 T_D=P_D\cap\ker d_{(3)},
\]
\[
 U_D=T_D\cap\ker W_{\rm odd},\qquad
 V_D=U_D\cap\ker T_{\rm odd}.
\]
Here \(D\) is the joint rational Kummer kernel inside \(M\),
\(\Phi_p\) is the free pro-\(p\) action, and
\(T,W\) are the corrected cubic and harmonic CM homomorphisms.
For a set of primes, its primary idempotent in
\(\widehat{\mathbf Z}\) has component one at those primes
and zero at the others; \(\varepsilon_p\) denotes a single
primary idempotent. All profinite powers below are taken
in the closed procyclic subgroup generated by their argument.

For a coefficient prime \(\ell\), write
\[
 \mathcal N_{3,\ell}
  =\{z\in(\mathbf Z_\ell[\omega])^\times:z\bar z=1\},
 \qquad \omega^2+\omega+1=0,
\]
and let \(\mathcal N_{3,\ell,\rm odd}\) be its maximal
pro-odd factor. Put
\begin{equation}\label{eq:exact-joint-odd-1}
\mathcal A_4=N_4(4)_{\rm odd},\qquad
 \mathcal C_3=
 \prod_{\ell\ge5}\bigl(\mathcal N_{3,\ell,\rm odd}\bigr)^3.
\end{equation}
We regard \(\mathcal C_3\) as a subgroup of \(N_3(4)_{\rm odd}\)
by putting the identity at coefficient primes two and three.
At two the level-four subgroup already has trivial odd factor.
The identity at three is an additional restriction.

\begin{theorem}\label{thm:exact-joint-odd-cm-image}
The full and arithmetic joint images on \(T_D\) agree:
\begin{equation}\label{eq:exact-joint-odd-2}
\begin{split}
 (T_{\rm odd},W_{\rm odd})(T_D)
   &=(T_{\rm odd},W_{\rm odd})(T_D\cap H)\\
   &=\mathcal C_3\times\mathcal A_4.
 \end{split}
\end{equation}
Moreover
\begin{equation}\label{eq:exact-joint-odd-3}
T_{\rm odd}(U_D)=T_{\rm odd}(U_D\cap H)=\mathcal C_3.
\end{equation}
Consequently
\begin{equation}\label{eq:exact-joint-odd-4}
T_D/V_D\cong\mathcal C_3\times\mathcal A_4,\qquad
 U_D/V_D\cong\mathcal C_3,
\end{equation}
and
\begin{equation}\label{eq:exact-joint-odd-5}
T_D=V_D(T_D\cap H),\qquad
 U_D=V_D(U_D\cap H).
\end{equation}
These are internal factorizations of the indicated groups.
\end{theorem}

\subsection{The upper bound, including the wild coefficient prime}

\begin{lemma}\label{lem:wild-three-primary-kernel}
For every \(g\in P_D\), the entire coefficient-three cubic
frame is the identity:
\[
                              T_{g,3}=1.
\]
\end{lemma}

\begin{proof}
The action \(\Phi_3(g)\) is the identity.
Hence its beta multiplier in the pro-three metabelian
quotient is the constant series one in
\(\mathbf Z_3[[X-1,Y-1]]\).
Specialization \(X=Y=\omega\) is continuous: in
\(\mathbf Z_3[\omega]\), the element \(\omega-1\) is
topologically nilpotent. Equivalently the corresponding
character of the abelianization has three-power order,
so it factors through the pro-three abelianization.
Thus the coefficient-three value of \(B_g(\omega,\omega)\)
is one.

Apply the established full integral cubic comparison,
Proposition~\ref{prop:lr-integral-cubic-beta}.
It identifies \(T_g\) with \(B_g(\omega,\omega)\) or its
conjugate on \(D\); the sign is one there.
Since \(P_D\subset D\), the asserted identity follows.
This uses the integral comparison at the ramified coefficient
prime, not a reduction modulo three.
\end{proof}

For \(g\in T_D\), the definition of \(d_{(3)}\), together with
\(\rho_\ell^I(g)=0\), gives \(c_{\ell,3}(g)=0\) for every
\(\ell\ge5\).
The defining exponent
\[
 \omega^{c_{\ell,3}(g)}
 =B_{g,\ell}(\omega,\omega)^{
                  (\ell-\chi_{-3}(\ell))/3}
 \quad\text{after reduction modulo }\ell
\]
and the integral cubic comparison say precisely that
\(T_{g,\ell}\) is a cube in \(\mathcal N_{3,\ell}\).
At these primes cubing is an automorphism on the principal
pro-\(\ell\) subgroup; the only condition is on the
three-primary part of the finite residue torus.
Lemma~\ref{lem:wild-three-primary-kernel} supplies the
coefficient-three identity, and the coefficient-two odd
factor is trivial. Therefore
\begin{equation}\label{eq:exact-joint-odd-6}
(T_{\rm odd},W_{\rm odd})(T_D)
        \subset\mathcal C_3\times\mathcal A_4.
\end{equation}

\subsection{Local arithmetic inertia in the cubic CM representation}

We use Tate representations for the arithmetic
constructions below. On \(D\cap H\) the established
comparison identifies their odd projections with the
corrected frames. The central two-primary correction in
the cubic comparison has no effect on these projections.

\begin{lemma}\label{lem:local-cubic-cm-inertia}
Let \(\ell\ge5\), let
\(e_\ell=v_3(\ell-\chi_{-3}(\ell))\), and let
\(Z_\ell\) be the cyclic three-Sylow subgroup of
\(\mathcal N_{3,\ell}\), of order \(3^{e_\ell}\).
There exists an arithmetic element
\[
 x_\ell\in M\cap H\cap\bigcap_p\ker\Phi_p
\]
whose uncorrected cubic Tate value is supported only at \(\ell\)
and generates \(Z_\ell\).
\end{lemma}

\begin{proof}
Write \(K=\mathbf Q(\omega)\), \(O=\mathbf Z[\omega]\), and
\(\pi=1-\omega\).
The earlier arithmetic CM computation gives
\begin{equation}\label{eq:exact-joint-odd-7}
\rho_3(G_K)=
 \{u\in\widehat O^\times:u\equiv1\pmod{2\pi}\}.
\end{equation}
The reduction of \(\mu_6\) onto
\((O/(2\pi))^\times\) is an isomorphism.
Thus \eqref{eq:exact-joint-odd-7} is the unique set of representatives of
\(\widehat O^\times/\mu_6\).
CM reciprocity
\cite[Theorem~12.27, p.~124]{MilneMF110} identifies the
inertia of a prime \(\mathfrak l\) over \(\ell\) with the
local units supported at \(\mathfrak l\), modulo \(\mu_6\).
Because \(\mathfrak l\nmid6\), those supported units already
satisfy the normalization in \eqref{eq:exact-joint-odd-7}.
Consequently their Tate image is the same supported
unit group, up to inversion and the choice of the CM
embedding. In particular its three-primary finite image
is full.

Choose \(\tau\) in this arithmetic inertia with three-primary
Tate value generating that local three-Sylow group, and
replace it by \(\tau^{\varepsilon_3}\).
If \(\ell\) splits in \(K\), its Tate value is, up to orientation,
\((\zeta,1)\), where \(\zeta\) generates the three-Sylow
subgroup of \(\mathbf Z_\ell^\times\).
If \(\ell\) is inert, its value is a generator of the
three-Sylow subgroup in the unramified quadratic unit group.
This latter subgroup has norm one, since \(3\nmid\ell-1\).

Let \(c\) be complex conjugation and set
\[
                         x_\ell=[c,\tau]^{\varepsilon_3}.
\]
The commutator has full multiplier one.
At every coefficient prime \(p\ne\ell\), its free pro-\(p\)
action is trivial because arithmetic inertia at \(\ell\)
acts trivially there
\cite[\S8.5, p.~151]{HainMatsumoto2003}.
At \(p=\ell\), its normalized image is pro-\(\ell\),
so the subsequent three-primary power kills that image.
Thus \(x_\ell\) has every primary action trivial.

Its cubic Tate value is the ratio \(\bar u/u\) of the
chosen local unit.
At a split prime this is \((\zeta^{-1},\zeta)\).
At an inert prime it is \(u^{-2}\), because the
three-Sylow subgroup has norm one.
In both cases it generates \(Z_\ell\), since two is
prime to three. It is supported only at \(\ell\), and
the power \(\varepsilon_3\) leaves it unchanged.
\end{proof}

\subsection{Realizing every permitted cubic component}

We will repeatedly use two elementary facts.
The image of \(\ker\lambda_p\) under \(\Phi_p\) is pro-\(p\):
it acts trivially on the Frattini quotient of the free
pro-\(p\) group, and the kernel of that finite linear action
on automorphisms is pro-\(p\).
Also, for \(a\in M\), every rational Kummer cocycle satisfies
\[
 \rho_q^I(hah^{-1})=\lambda(h)\rho_q^I(a),\qquad
 \rho_q^I([h,a])=(\lambda(h)-1)\rho_q^I(a).
\]
Thus complementary primary orders kill the corresponding
normalized free-group actions, and the displayed formulas
give exact cancellation of the full adelic Kummer values.

Choose once and for all an arithmetic pro-two element
\(h_3\) whose full multiplier is \(-1\) at three and one
at every other prime. It fixes \(\mathbf Q(i)\) and
conjugates \(\mathbf Q(\sqrt{-3})\).

First consider a prescribed cubic odd value in the product of the Sylow subgroups for odd primes
different from three.
The arithmetic corrected cubic image supplies a
multiplier-one lift \(m\).
Take its primary power \(a=m^\varepsilon\), where
\(\varepsilon\) is one at primes other than two and three
and zero at those two primes.
The uncorrected cubic Tate value of \(a\) is the prescribed
value; the central two-primary correction has disappeared.

The commutator \([h_3,a]\) has every primary action trivial.
At two and three use the primary order of \(a\); at all
other primes use that \(h_3\) is pro-two with local
multiplier one.
Its rational Kummer values vanish, since those of \(a\)
are supported where the multiplier of \(h_3\) is one.
It is arithmetic, so its cubic defect vanishes.
Its harmonic CM value is one, because \(h_3\) fixes
\(\mathbf Q(i)\). Its cubic odd value is the inverse
square of the prescribed value.
Since squaring is an automorphism on an odd group, this
realizes the entire part of \(\mathcal C_3\) of primary
order prime to three inside \(U_D\cap H\).

It remains to realize \(Z_\ell^3\) for each \(\ell\ge5\).
Let \(x_\ell\) be as in
Lemma~\ref{lem:local-cubic-cm-inertia}, with cubic value
\(z_\ell\) generating \(Z_\ell\).
Choose \(h_4\in H\) with multiplier four at three and one
at every other prime, and put
\[
 y_\ell=h_4x_\ell h_4^{-1}x_\ell^{-\lambda(h_4)}.
\]
Every primary action remains trivial.
All rational Kummer values cancel by their multiplier-one
conjugation law, so \(y_\ell\in P_D\cap H\).
The element \(h_4\) fixes the cubic CM field, and \(x_\ell\)
has a three-primary Tate value. Therefore
\[
                       \rho_3(y_\ell)=z_\ell^{-3}.
\]
Now set
\[
                            w_\ell=[h_3,y_\ell].
\]
It still belongs to \(P_D\cap H\), hence to \(T_D\cap H\).
The harmonic value is one because \(h_3\) fixes
\(\mathbf Q(i)\).
The cubic odd value is \(z_\ell^6\), which generates
\(Z_\ell^3\). Thus \(w_\ell\in U_D\cap H\) realizes the
full permitted three-primary subgroup at the single
coefficient prime \(\ell\), with identity at the others.

The image of the compact group \(U_D\cap H\) is closed.
It contains every such finite coordinate axis, their
finite products, and therefore the full product of the
groups \(Z_\ell^3\). Combining this with the preceding
prime-to-three construction gives
\begin{equation}\label{eq:exact-joint-odd-8}
T_{\rm odd}(U_D\cap H)=\mathcal C_3.
\end{equation}
No convergence of an ordered infinite noncommuting product
has been assumed.

\subsection{Separating the harmonic component}

For any \(v\in\mathcal A_4\), the established harmonic
arithmetic image supplies \(m\in M\cap H\) with that odd
harmonic value.
Take \(a=m^{1-\varepsilon_2}\).
Choose the arithmetic pro-two element \(h_2\) with
multiplier \(-1\) at two and one at every odd prime.
The commutator \([h_2,a]\) belongs to \(T_D\cap H\) by
the argument for the full odd harmonic image:
its primary actions and all rational Kummer values vanish.
It has harmonic odd value \(v^{-2}\).

In addition, \(h_2\) fixes the cubic CM field, since its
multiplier at three is one.
Its uncorrected cubic Tate action therefore commutes with that
of \(a\), and its cubic odd commutator value is one.
It follows that the joint arithmetic image contains
\[
                         \{1\}\times\mathcal A_4.
\]
Together with \eqref{eq:exact-joint-odd-8}, it contains
\(\mathcal C_3\times\mathcal A_4\).
The upper bound \eqref{eq:exact-joint-odd-6} proves \eqref{eq:exact-joint-odd-2} and \eqref{eq:exact-joint-odd-3}.

The kernel of the joint map on \(T_D\) is \(V_D\).
Compactness gives \eqref{eq:exact-joint-odd-4}.
For each element of \(T_D\), choose an arithmetic member
of \(T_D\) with the same pair of odd CM values and divide
by it. This proves the first equality in \eqref{eq:exact-joint-odd-5}.
The same argument on \(U_D\), using \eqref{eq:exact-joint-odd-8}, proves the second.

\begin{remark}[Internal factorization and relative coinvariants]
The quotients in \eqref{eq:exact-joint-odd-4} are pro-odd, so the established
maximal pro-two relative quotient and original \(X_2\)
restriction kernel remain unchanged.
The internal factorizations in \eqref{eq:exact-joint-odd-5} do not, by themselves,
identify the full relative coinvariant quotients.

For a precise counterexample, let \(p\) be odd and let
\(G=\langle a,b,c\rangle\) be the order-\(p^3\) Heisenberg
group, with \(c\) central, \([a,b]=c\), and each generator
of order \(p\). Put
\[
 H=\langle b\rangle,\qquad T=\langle b,c\rangle,\qquad
 V=\langle c\rangle.
\]
Both \(T,V\) are \(G\)-normal and \(T=V(T\cap H)\).
Nevertheless
\[
 V/[G,V](V\cap H)\cong C_p,\qquad
 T/[G,T](T\cap H)=1.
\]
The character of \(V\) sending \(c\) to one extends to \(T\)
by sending \(b\) to zero, but this extension is not
\(G\)-invariant: conjugation by \(a\) sends \(b\) to \(bc\).
Thus cross-conjugation of \(T\cap H\) by \(G\) is a real
additional obstruction. No full relative coinvariant
identification is inferred from \eqref{eq:exact-joint-odd-5}.
\end{remark}
